% 解説プリントのコメント（saiten が書く。./tools/kaisetsu.ps1 がこのファイルからプリントを作る）
% 2026-11-22 のスキャン（p1）の採点から作った
%% matome
\begin{daiji}
  \item 分散関係の定数項（静止エネルギー $mc^{2}/\hbar$）は波動関数全体に共通の位相 $\ee^{-\ii mc^{2}t/\hbar}$ を掛けるだけで $\abs{\psi}^{2}$ を変えないので，観測される粒子の運動（包絡の速さ $v_{\mathrm{g}}$）には効かない。$v_{\mathrm{p}}$ はエネルギーの原点の取り方で変わる量で，粒子の速さではない（問題1 (3)）。
\end{daiji}
\begin{maemo}
  \item 位相速度の扱い：第2章 問題5（2026-10-29）で位相速度を $\hbar k/m$ と取り違えた。位相速度は $\omega/k$ で，エネルギーの原点をずらすと変わる。物理的な速さとして使えるのは $\dd\omega/\dd k$ の方である，と毎回確かめる。
\end{maemo}
%% mondai 1
\begin{kotaeran}
  \machigai{3}{○}{記述} 展開 $\omega\approx c\mu+ck^{2}/(2\mu)=mc^{2}/\hbar+\hbar k^{2}/(2m)$，$v_{\mathrm{g}}\approx\hbar k/m$，$v_{\mathrm{p}}\approx mc^{2}/(\hbar k)+\hbar k/(2m)$ と，非相対論の値と比べて $v_{\mathrm{g}}$ が一致することは正しい。理由が「$v_{\mathrm{p}}$ は静止エネルギー分ずれる」だけで，問われている「一致する理由」の中心が書かれていない。(i) 定数項 $mc^{2}/\hbar$ は $k$ で微分すると消えるので $v_{\mathrm{g}}$ には効かず，$\omega/k$ には $mc^{2}/(\hbar k)$ として残る。(ii) 定数項は全体に共通の位相を掛けるだけで $\abs{\psi}^{2}$ を変えないので，観測される運動は変わらない。この2点まで書く。
\end{kotaeran}
\begin{rikai}
  \item $\omega\to\omega+\omega_{0}$ と定数をずらすと $\psi\to\psi\,\ee^{-\ii\omega_{0}t}$ となり，$\abs{\psi}^{2}$ は変わらない。観測にかかるのは $\omega$ の $k$ による変化だけで，$\dd\omega/\dd k$ はそれだけで決まる。
  \item 「どちらが一致するか，その理由」の問いには，式の上の理由（微分で定数が消える）と物理的な理由（共通の位相は観測されない）の両方を書く。
\end{rikai}
